% =============================================================================
%  4. Estudio Financiero
% =============================================================================
\section{Estudio Financiero}

\begin{quote}
\emph{Cifras en USD. Modelo a 5 años, escenario base. Las cifras son estimaciones para
fines académicos del proyecto.}
\end{quote}

\subsection{Componentes de pre-inversión}

\begin{table}[ht]
  \centering
  \begin{tabularx}{\linewidth}{Y r}
    \toprule
    \textbf{Componente} & \textbf{Costo (USD)} \\
    \midrule
    Estudio de mercado y validación & 1\,500 \\
    Constitución legal de la empresa y marca & 2\,000 \\
    Diseño de arquitectura y prototipo (MVP inicial) & 8\,000 \\
    Asesoría legal (términos, protección de datos, IP) & 1\,500 \\
    \midrule
    \textbf{Total pre-inversión} & \textbf{13\,000} \\
    \bottomrule
  \end{tabularx}
\end{table}

\subsection{Inversiones requeridas (Año 0)}

\begin{table}[ht]
  \centering
  \begin{tabularx}{\linewidth}{Y r}
    \toprule
    \textbf{Inversión} & \textbf{Costo (USD)} \\
    \midrule
    Desarrollo del producto hasta lanzamiento comercial & 45\,000 \\
    Equipos de trabajo (laptops, periféricos) & 6\,000 \\
    Configuración de infraestructura cloud inicial & 3\,000 \\
    Licencias y herramientas de desarrollo & 2\,000 \\
    Identidad de marca y sitio web & 2\,500 \\
    Capital de trabajo inicial & 11\,500 \\
    \midrule
    \textbf{Total inversión inicial} & \textbf{70\,000} \\
    \bottomrule
  \end{tabularx}
\end{table}

\noindent\textbf{Inversión total del proyecto (pre-inversión + inversión):
USD 83\,000.}

\subsection{Costo de cada componente para producir (anual, estabilizado)}

\subsubsection{Mercadeo y/o divulgación}
Marketing digital, contenido y webinars (de bajo costo, enfoque orgánico):
\textbf{USD 6\,000/año} (año 1), creciendo con las ventas.

\subsubsection{Costo de reuniones y prospección de clientes}
Demos, viáticos y herramienta de CRM básica: \textbf{USD 4\,000/año} (año 1).

\subsubsection{Costos de producción (operación del servicio)}
Infraestructura cloud, feeds de vulnerabilidades y cómputo escalable:
\textbf{USD 4\,000/año} (año 1), variable según número de clientes
($\sim$USD 12/seat/año).

\subsubsection{Costos de mantenimiento}
Actualizaciones, soporte técnico y corrección de errores: \textbf{USD 6\,000/año}
(año 1).

\subsection{Gastos del proceso: fijos y variables}

\begin{table}[ht]
  \centering
  \begin{tabularx}{\linewidth}{l Y r}
    \toprule
    \textbf{Tipo} & \textbf{Concepto} & \textbf{Monto anual (USD)} \\
    \midrule
    Fijos     & Salarios equipo base (3 personas, fundadores con salario modesto) & 60\,000 \\
    Fijos     & Coworking / oficina pequeña & 4\,800 \\
    Fijos     & Servicios, contabilidad, seguros & 5\,000 \\
    Fijos     & Licencias y herramientas & 3\,000 \\
    Variables & Infraestructura cloud (escala con clientes) & $\sim$12 / seat / año \\
    Variables & Comisiones de venta y marketplaces & $\sim$5\,\% de ingresos \\
    Variables & Soporte adicional por cliente & $\sim$8 / seat / año \\
    \bottomrule
  \end{tabularx}
\end{table}

\subsection{Análisis económico financiero}

\subsubsection{Egresos}
\textbf{Inversiones:} USD 83\,000 (año 0).

\noindent\textbf{Costos + Gastos operativos por año (USD):}

\begin{table}[ht]
  \centering
  \begin{tabular}{c r r r r}
    \toprule
    \textbf{Año} & \textbf{Salarios y fijos} & \textbf{Mkt. + ventas} & \textbf{Prod. + mant. (var.)} & \textbf{Total egresos op.} \\
    \midrule
    1 & 70\,000  & 10\,000 & 6\,000  & 86\,000  \\
    2 & 85\,000  & 12\,000 & 10\,000 & 107\,000 \\
    3 & 100\,000 & 18\,000 & 16\,000 & 134\,000 \\
    4 & 120\,000 & 25\,000 & 24\,000 & 169\,000 \\
    5 & 145\,000 & 32\,000 & 34\,000 & 211\,000 \\
    \bottomrule
  \end{tabular}
\end{table}

\subsubsection{Ingresos}
\textbf{a) Venta de productos (suscripciones):} ingreso principal $=$ seats activos
$\times$ precio promedio (USD 28/mes $\times$ 12).

\begin{table}[ht]
  \centering
  \begin{tabular}{c r r}
    \toprule
    \textbf{Año} & \textbf{Seats activos} & \textbf{Ingreso anual suscripciones (USD)} \\
    \midrule
    1 & 40  & 13\,440  \\
    2 & 99  & 33\,264  \\
    3 & 192 & 64\,512  \\
    4 & 325 & 109\,200 \\
    5 & 532 & 178\,752 \\
    \bottomrule
  \end{tabular}
\end{table}

\noindent\textbf{b) Venta de activos:} no aplica de forma relevante (SaaS sin activos
físicos significativos); eventual venta de equipo depreciado al año 5
($\sim$USD 3\,000).

\noindent\textbf{c) Venta de desechos o subproductos:} como «subproducto», la empresa
puede ofrecer \textbf{servicios de consultoría y reportes de inteligencia de amenazas}
derivados de los datos agregados y anonimizados, y capacitaciones. Ingreso estimado:

\begin{table}[ht]
  \centering
  \begin{tabular}{c r}
    \toprule
    \textbf{Año} & \textbf{Ingreso subproductos/servicios (USD)} \\
    \midrule
    1 & 2\,000  \\
    2 & 5\,000  \\
    3 & 10\,000 \\
    4 & 16\,000 \\
    5 & 24\,000 \\
    \bottomrule
  \end{tabular}
\end{table}

\noindent\textbf{Ingresos totales por año (USD):}

\begin{table}[ht]
  \centering
  \begin{tabular}{c r r r}
    \toprule
    \textbf{Año} & \textbf{Suscripciones} & \textbf{Servicios/subproductos} & \textbf{Total ingresos} \\
    \midrule
    1 & 13\,440  & 2\,000           & 15\,440  \\
    2 & 33\,264  & 5\,000           & 38\,264  \\
    3 & 64\,512  & 10\,000          & 74\,512  \\
    4 & 109\,200 & 16\,000          & 125\,200 \\
    5 & 178\,752 & 24\,000 + 2\,000 & 204\,752 \\
    \bottomrule
  \end{tabular}
\end{table}

\subsubsection*{Flujo de caja del proyecto (USD)}

\begin{table}[ht]
  \centering
  \begin{tabular}{c r r r r}
    \toprule
    \textbf{Año} & \textbf{Ingresos} & \textbf{Egresos op.} & \textbf{Inversión} & \textbf{Flujo neto} \\
    \midrule
    0 & 0        & 0        & 83\,000 & \textbf{$-83\,000$} \\
    1 & 15\,440  & 86\,000  & 0       & \textbf{$-70\,560$} \\
    2 & 38\,264  & 107\,000 & 0       & \textbf{$-68\,736$} \\
    3 & 74\,512  & 134\,000 & 0       & \textbf{$-59\,488$} \\
    4 & 125\,200 & 169\,000 & 0       & \textbf{$-43\,800$} \\
    5 & 204\,752 & 211\,000 & 0       & \textbf{$-6\,248$}  \\
    \bottomrule
  \end{tabular}
\end{table}

\subsection{Criterio de decisión: VAN y TIR}
Se utiliza una \textbf{tasa de descuento del 15\,\%} (refleja el alto riesgo de una
startup tecnológica y el costo de oportunidad del capital).

\noindent\textbf{Cálculo del Valor Actual Neto (VAN):}

\begin{table}[ht]
  \centering
  \begin{tabular}{c r r r}
    \toprule
    \textbf{Año} & \textbf{Flujo neto (USD)} & \textbf{Factor $(1.15)^t$} & \textbf{Valor presente (USD)} \\
    \midrule
    0 & $-83\,000$ & 1.0000 & $-83\,000$ \\
    1 & $-70\,560$ & 1.1500 & $-61\,357$ \\
    2 & $-68\,736$ & 1.3225 & $-51\,973$ \\
    3 & $-59\,488$ & 1.5209 & $-39\,114$ \\
    4 & $-43\,800$ & 1.7490 & $-25\,043$ \\
    5 & $-6\,248$  & 2.0114 & $-3\,106$  \\
    \bottomrule
  \end{tabular}
\end{table}

\noindent\textbf{VAN (al 15\,\%) $\approx -$USD 263\,593}

\begin{quote}
El VAN es \textbf{negativo} en el horizonte de 5 años. Esto es característico de los
proyectos SaaS pequeños: la inversión y los costos fijos iniciales pesan más que los
ingresos tempranos, y la rentabilidad llega después del periodo evaluado. Lo positivo
es que el flujo neto mejora cada año y el año 5 casi alcanza el punto de equilibrio
($-6\,248$), lo que anticipa flujos positivos a partir del año 6.
\end{quote}

\noindent\textbf{Tasa Interna de Retorno (TIR):} dado que todos los flujos del proyecto
son negativos durante los 5 años (flujo acumulado $\approx -$USD 331\,832),
\textbf{no existe una TIR positiva dentro del horizonte de 5 años}; el proyecto aún no
recupera la inversión en ese plazo.

\subsubsection{Justificación y recomendación}
Con los supuestos del escenario base, en un horizonte de \textbf{5 años el proyecto NO
es viable} según el VAN (negativo) ni la TIR (no alcanza retorno). Sin embargo, hay
matices importantes para la recomendación:

\begin{enumerate}
  \item \textbf{Tendencia positiva clara:} el flujo neto anual mejora cada año, pasando
        de $-70\,560$ (año 1) a $-6\,248$ (año 5). El negocio escala de forma sana y se
        acerca al punto de equilibrio al final del periodo.
  \item \textbf{Naturaleza del modelo SaaS:} estos negocios típicamente alcanzan el
        punto de equilibrio entre los años 5 y 7. Un horizonte de evaluación de 5 años
        es corto para este tipo de inversión, sobre todo en un escenario de crecimiento
        conservador.
  \item \textbf{Sensibilidad a los supuestos:} el resultado depende fuertemente de los
        salarios (el mayor egreso) y del ritmo de captación de clientes. Una captación
        algo más rápida o costos fijos aún más ajustados harían que el proyecto cruzara
        a flujos positivos antes del año 5.
\end{enumerate}

\noindent\textbf{Recomendación:} Bajo los supuestos actuales y un horizonte de
\textbf{5 años, el proyecto aún no es rentable} según el VAN (negativo) ni la TIR (no
se alcanza retorno). Sin embargo, dado que el flujo prácticamente alcanza el equilibrio
en el año 5 y la tendencia es claramente ascendente, \textbf{el proyecto es prometedor
a mediano plazo}. Se recomienda:
\begin{itemize}
  \item[(a)] \textbf{Extender la evaluación a 7--8 años}, plazo realista para un SaaS
        bootstrapped, donde el VAN probablemente se vuelva positivo dada la curva de
        crecimiento.
  \item[(b)] \textbf{Mantener la estructura de costos ajustada} (equipo pequeño,
        infraestructura escalonada) durante los primeros años.
  \item[(c)] \textbf{Buscar acelerar moderadamente la captación de clientes} para
        adelantar el punto de equilibrio sin comprometer la calidad del servicio.
\end{itemize}

En conclusión, el proyecto \textbf{no se aprueba en un horizonte estricto de 5 años por
VAN/TIR}, pero es una inversión \textbf{estratégicamente atractiva y de bajo monto
inicial} (USD 83\,000) que merece reevaluarse con un horizonte más largo, considerando
el fuerte crecimiento del mercado (CAGR de dos dígitos) y la demanda regional
insatisfecha.
